% Reciprocal proofs for the PA/Puel source-specific continuation.
\section{PA norm/indexing proofs and the distinct embedded-left-side case}

\begin{annexitem}{annex:p12-klop1987-pa-reannotation-failure}{thm:p12-klop1987-pa-reannotation-failure}{An actual inner PA step that enlarges its ancestor's index}
\begin{proof}
Use the first actual tutorial rule at the inner parallel node of
\(t=(a\parallel b)\lhd c\). Syntax-node lengths of the three
action leaves are1. Its inner parallel subtree has length3;
the outer left-merge's initial argument-sum index is therefore4.
After the step the new left argument is
\((a\lhd b)+(b\lhd a)\), of length \(1+3+3=7\).
Its reannotated outer index is8. Each new inner left-merge has
index2. The old tree contains outer \(L_4\), inner \(P_2\)
and action leaves; the new tree's outer head is \(L_8\).
For example numerical labels \(P_n=2n+1,L_n=2n\), with
base operators/actions below \(L_2\), change the largest root
label from8 to16. More generally the source's interleaving implies
\(L_8>L_4>P_2\) by transitivity, so any numerical realization
of that precedence gives the new root a larger number than all
numbers in the old tree.

For the primary report's additive atomic-length norm applied to
the tutorial's two-branch PA rule, every action again has value1
but the parallel subtree has value2 rather than3. The outer index
is3. Its new left argument has value
\((1+1)+(1+1)=4\), and the outer index is5. Now the old heads
are \(L_3,P_2\), whereas \(L_5>L_3>P_2\).
This is the same contextual failure for that different norm.

The primary's actual ACP rule CM1 has a third branch \(a|b\).
For that rule the same original outer index is3, but the new left
argument is \((a\lhd b)+(b\lhd a)+(a|b)\). Its primary
additive norm is \(2+2+2=6\), whatever the choice parentheses,
so the outer index is7. The new head \(L_7\) is again strictly
above all old heads \(L_3,P_2,a,b,c\) in the stated precedence.
This pins the same missing contextual lift against the actual
printed primary rule, rather than quietly replacing its ACP
merge equation by the tutorial's PA variant.

Every marker operation leaves numerical labels unchanged, deletes
nodes, copies existing nodes, or introduces a strictly smaller
head. Context closure cannot introduce a larger number either.
Induction on the length of a marker reduction therefore bounds
all its labels by the maximum label of its input. Thus none of the
freshly reannotated outputs above can be reached from its input by
the required marker reduction. In LPO the same example also fails:
its new outer head is above every head in the old term, and neither
the old root nor any of its descendants can satisfy the head/lex
clause against it; the subterm clause recursively has no witness.

This argument treats an actual contextual instance of a printed
PA rule, not a hypothetical unrelated duplicating system. It does
not refute the source's root-level typical comparisons or the
termination statement. Retaining old tags as budgets is a different
possible lift; it requires the admissibility invariant proved below.
\end{proof}
\end{annexitem}

\begin{annexitem}{annex:p12-klop1987-pa-norm}{thm:p12-klop1987-pa-norm}{All norm equations, positive substitution values and the common context}
\begin{proof}
Structural recursion defines \(\nu\), since terms are finite.
The base values are positive integers. Maximum and addition preserve
positivity and dominate each operand, so induction gives positivity
and the bound for every immediate child, hence every subterm.

For a substitution \(\sigma\), assign the variable \(x\) the
positive integer \(\nu(\sigma(x))\). The expression formed
from a term by replacing its variables by those integers, actions
by1, choice by maximum and every other binary operator by addition
evaluates to \(\nu(t\sigma)\). Induct on \(t\): the variable
and action cases are precisely this assignment; each constructor
combines the two induction values with its declared operation.
Thus all following rule equations hold for arbitrary positive
substitution values, not just when every substituted variable is
still a one-node leaf.

Put \(X=\nu(x\sigma),Y=\nu(y\sigma),Z=\nu(z\sigma)\).
For parallel expansion the left value is \(X+Y\), and the
right value is \(\max(X+Y,Y+X)=X+Y\). For left-merge
distribution the two values are
\(\max(X,Y)+Z\) and \(\max(X+Z,Y+Z)\), equal because
addition by the fixed \(Z\) preserves order on integers.
For consuming an atomic prefix they are
\((1+X)+Y\) and \(1+(X+Y)\), equal by associativity of
addition; atomicity fixes the first operand's norm to1.
For atomic left-merge the two values are \(1+X\).
Idempotence gives \(\max(X,X)=X\). Sequential distribution
has the same max/add identity as left-merge distribution.
Sequential associativity gives \((X+Y)+Z=X+(Y+Z)\).
These are all seven rules. Every action retains positive unit
weight; no arbitrary zero or negative variable/action weights
are part of this construction.

If two terms have the same norm, putting them in the same immediate
parent position gives the same parent norm, because the unchanged
other child's norm is the same and the same maximum or sum is
applied. Induction along the ancestor path proves equality of the
whole context norms. At a parallel or left-merge ancestor the
argument-sum index is therefore unchanged as well.

More explicitly, for a one-hole ordinary context \(C\) and a
positive hole value \(w\), form its annotated context recursively:
the hole remains a hole with scalar value \(w\); fixed siblings
are annotated by \(E\); every parent's scalar value is computed
by the same maximum/sum recursion, and each merge head receives
the sum of its two scalar argument values. This is a well-defined
finite recursion. Induction on the context depth proves
\(E(C[v])=\widehat C_w[E(v)]\) whenever \(\nu(v)=w\).
For a rule's two instantiated endpoints the same \(w\) applies,
so they have exactly the same annotated context \(\widehat C_w\).
This proves the claimed contextual lift without asserting the
false general identity \(E(t\sigma)=E(t)\sigma\).
Indeed \(E(x\parallel y)=P_2(x,y)\); substituting
\(a\cdot a\) for \(x\) gives
\(E((a\cdot a)\parallel y)=P_3(a\cdot a,y)\), while
ordinary substitution into the initially annotated expression leaves
its head \(P_2\). The instance-first and fixed-hole-norm arguments
above avoid that concrete mismatch.
\end{proof}
\end{annexitem}

\begin{annexitem}{annex:p12-klop1987-pa-lpo-lift}{thm:p12-klop1987-pa-lpo-lift}{All seven source-instance comparisons and their required child domination}
\begin{proof}
Use the recursive LPO definition and its proved transitivity,
proper-subterm and context properties. Write \(p,q,r\) for
the annotations of the actual finite terms instantiating variables,
and \(X,Y,Z\) for their erasures' positive norms. This proves
each instance with its actual indices; it is not a check of just
the schematic variables with norm1.

For parallel expansion let \(n=X+Y\) and
\(S=P_n(p,q)\). It dominates \(p,q\) by the subterm property.
Since \(P_n>L_n\), the LPO head clause gives both
\(S>L_n(p,q)\) and \(S>L_n(q,p)\). The same clause with
\(P_n>+\), now dominating both these target children, gives
\(S>L_n(p,q)+L_n(q,p)\).

For left-merge distribution put
\(n=\max(X,Y)+Z\), \(i=X+Z\), \(j=Y+Z\) and
\(S=L_n(p+q,r)\). Then \(i,j\leq n\), and \(S\)
dominates \(p,q,r\), including the first two through their
proper occurrences inside \(p+q\). To compare with \(L_i(p,r)\),
if \(i<n\) use strict precedence \(L_n>L_i\) and domination
of \(p,r\). If \(i=n\), the equal-head tuple comparison
has first coordinate \(p+q>p\) and unchanged second coordinate
\(r\); the same domination of both target arguments makes
the full LPO clause valid. The comparison with \(L_j(q,r)\)
uses exactly those two cases with \(p+q>q\). Finally
\(L_n>+\) and domination of both whole branch terms gives
the comparison with their sum. Neither branch was assumed to
have the other's norm.

For consuming a prefix take \(n=1+X+Y\geq3\) and
\(S=L_n(a\cdot p,q)\). It dominates \(a,p,q\) by proper
subterms and transitivity. The source interleaving gives
\(L_n>P_{n-1}\); thus \(S>P_{n-1}(p,q)\) by the head
clause with both target children dominated. Since \(L_n>\cdot\)
and it dominates both \(a\) and that new parallel term, the
head clause gives \(S>a\cdot P_{n-1}(p,q)\).
The right parallel index is exactly \(X+Y=n-1\).

For atomic left-merge, \(S=L_{1+X}(a,p)\) dominates \(a,p\).
Its head is above \(\cdot\), giving \(S>a\cdot p\).
For idempotence \(p+p>p\) is an immediate-subterm comparison.

For sequential distribution let \(S=(p+q)\cdot r\).
It dominates \(p,q,r\). At the equal product head, comparison
with \(p\cdot r\) uses first coordinate \(p+q>p\), and
comparison with \(q\cdot r\) uses \(p+q>q\); every target
argument has already been dominated. The product head is above
choice, so domination of those two whole products then yields
\(S>(p\cdot r)+(q\cdot r)\).

For sequential associativity let \(S=(p\cdot q)\cdot r\).
It dominates \(p,q,r\). Comparison with \(q\cdot r\)
uses first coordinate \(p\cdot q>q\), unchanged second
coordinate \(r\), and domination of \(q,r\); hence
\(S>q\cdot r\). To compare with \(p\cdot(q\cdot r)\),
use first coordinate \(p\cdot q>p\), together with the
now established domination of both target arguments
\(p,q\cdot r\). The equal-head LPO clause applies.

This exhausts the seven rules. Canonical annotation is deterministic,
so repeated instantiations of one variable, including the two
copies in the idempotence pattern, give literally the same
annotated subterm. The norm lemma supplies one common annotated
context around both endpoints of each ordinary contextual step.
LPO context closure then gives the strict whole-term comparison.
\end{proof}
\end{annexitem}

\begin{annexitem}{annex:p12-klop1987-pa-termination}{thm:p12-klop1987-pa-termination}{The actual finite alphabet for an entire PA run}
\begin{proof}
Let the ordinary starting term have norm \(N\). Norm invariance
gives that same \(N\) at every finite stage of any run. A merge
node has norm equal to its argument-sum index, and the subterm
bound gives \(n\leq N\). Positivity of both arguments gives
\(n\geq2\). Thus at most \(2(N-1)\) indexed merge symbols
can occur; if \(N=1\), none occurs. Syntax size may grow under
distribution or parallel expansion, but that does not change
this norm/index bound.

No rule introduces a new object variable or a new action: every
right-side variable already occurs on its left, and every displayed
action is the same captured atomic symbol retained from its left.
Hence only the finite action and variable sets in the initial
term can occur. Include \(+\) and \(\cdot\) even if absent
initially. Together with \(P_n,L_n\) for \(2\leq n\leq N\),
these give one finite ranked signature for the entire annotated
run, with one fixed left-to-right status and the restriction of
the declared strict precedence.

By the rule/context lift an infinite ordinary run would give an
infinite strict LPO descent over this one signature. The complete
finite-ranked structural/tree/well-foundedness argument in
Theorems~\ref{thm:p12-ro-lpo-structural},
\ref{thm:p12-ro-finite-tree} and \ref{thm:p12-ro-lpo-termination}
excludes it. Those results allow a partial precedence; totality
between different actions is not required. Their open-term case
treats the fixed finite variable set of the first term as additional
nullary symbols for the embedding argument. Thus the conclusion
holds for open as well as ground inputs, even with an infinite
ambient action supply. Restriction to \(R_3\subseteq R_7\)
gives its termination as well.

For the stated excluded equation boundary, directing
\(x+y\to y+x\) gives
\(x+y\to y+x\to x+y\to\cdots\) on two distinct object
variables. Thus treating the full choice equations as unrestricted
directed rewriting would invalidate an SN conclusion; none was
implicitly added to the proved systems.
\end{proof}
\end{annexitem}

\begin{annexitem}{annex:p12-klop1987-pa-budget-lift}{thm:p12-klop1987-pa-budget-lift}{Initial syntax lengths can be retained only with an actual admissibility invariant}
\begin{proof}
Let \(\ell(t)\) count every action, variable and operator node.
Induction shows \(\nu(t)\leq\ell(t)\): leaves have value1;
at a nonchoice binary symbol the sum of the two induction bounds
is at most1 plus their syntax lengths; at choice, maximum is at
most that same sum. Therefore initially tagging each merge by
\(\ell(p)+\ell(q)\) makes it admissible, with index at least2.

Suppose an admissible annotation is given and an ordinary source
step is selected. Every \(R_3\) left pattern is left-linear,
and its unindexed \(+\), product and atomic constant positions
are unchanged by annotation. Match the annotated version by
capturing the existing annotated variable-instance subterms.
No repeated-variable match between independently budgeted values
has to be invented. Use the lifted rule with the existing outer
index \(n\), not a fresh syntax-length tag.

For parallel expansion, admissibility of the input gives
\(n\geq\nu(p)+\nu(q)\). Both new \(L_n\)-nodes satisfy
this condition, including the swapped argument pair. Their sum's
erasure has the same norm as the old parallel erasure.
For distribution, the input gives
\(n\geq\max(\nu(p),\nu(q))+\nu(r)\), which is at least
each of \(\nu(p)+\nu(r)\) and \(\nu(q)+\nu(r)\).
Thus both inherited \(L_n\)-nodes are admissible. Their combined
erasure norm is unchanged. For consuming the atomic prefix,
\(n\geq1+\nu(p)+\nu(q)\), so \(n\geq3\) and
\(n-1\geq\nu(p)+\nu(q)\). The new \(P_{n-1}\) is
well-defined and admissible. Its surrounding product again has
the old erasure norm. Captured descendants keep their already
admissible tags in every retained or copied occurrence.

Every ordinary ancestor has unchanged erasure-argument norms by
the norm invariance lemma. Its old upper budget therefore remains
admissible. This proves, by induction along the context path, an
actual lift in the same unchanged-tag annotated context and
preservation of admissibility after every step.

For strict LPO descent, the parallel and prefix-consuming root
comparisons are exactly those already proved with \(P_n>L_n\)
and \(L_n>P_{n-1}\); their domination calculations use the
captured annotated terms and do not require equality of the
budget to their norm sum. Distribution now has \(L_n\) in
both branches, regardless of their unequal norms. Comparison
of \(L_n(p+q,r)\) with \(L_n(p,r)\) uses \(p+q>p\),
and with \(L_n(q,r)\) uses \(p+q>q\), while the source
dominates each target argument. Its higher head than choice
then dominates the whole right side. Context closure applies
because surrounding tags were retained literally.

No lifted rule introduces a larger index: expansion and distribution
inherit \(n\); consuming a prefix introduces \(n-1\).
Let \(B\) be the finite largest tag initially present. The entire
lift uses only indices2 through \(B\), fixed base operators and
the finite action/variable sets of the initial term. If no tag is
present, no \(R_3\)-redex can occur. Finite-ranked LPO
well-foundedness therefore excludes an infinite lift, and every
ordinary \(R_3\)-step was shown to have such a lift. This is
a complete route from the tutorial's initial syntax-length tags.
It is different from recomputing those tags after every duplicating
step, which the explicit counterexample refuted.

The left-linearity used here belongs to these three patterns,
not to every larger PA/ACP presentation. The canonical max/sum
proof separately handled the nonlinear idempotence rule; no
general nonlinear upper-budget coherence theorem was assumed.
\end{proof}
\end{annexitem}

\begin{annexitem}{annex:p12-klop1987-pa-puel-example}{thm:p12-klop1987-pa-puel-example}{A finite active-origin count for the actual isolated rule}
\begin{proof}
Call an \(f\)-occurrence active precisely when its immediate
argument has root \(s\). These are exactly the redexes of \(Q\).
In this one-rule system a subterm with root \(s,p,g\), an inert
constant or a variable cannot change that root; an \(f\)-root
can only remain \(f\) or change to \(g\), never to \(s\).
Consequently an initially inactive \(f\)-occurrence cannot
become active through a reduction inside its argument. An active
one remains active until its own root is contracted: steps below
its argument's \(s\)-root preserve that root.

In a contraction the active root \(f(s(x))\) is replaced by
\(g(s(x),f(p(s(x))))\). The new \(f\)-root is inactive
because its argument has root \(p\), and remains so, since
\(p\) has no rule. Every other surviving active occurrence
is within one of the two copies of the old \(x\)-instance.
Thus active occurrences are copied descendants of initially
active origins, and no new active origin is created.

List the initially active occurrences by their containment forest,
giving each its depth in that finite forest. Every proper active
descendant has strictly greater initial depth. Retain this assigned
origin depth when an occurrence is copied; it is not its changing
physical depth in the new syntax. Let \(c_i\) count current
active residuals of origin depth \(i\). There is one fixed
finite tuple \((c_0,\ldots,c_d)\), where \(d\) is the greatest
initial origin depth; an input with no active occurrence is already
normal.

Contracting an active origin of depth \(i\) removes one residual
from \(c_i\). Only proper active descendants in its captured
instance can be copied, and their origin depths are greater than
\(i\). Counts at earlier depths are unchanged. Other active
ancestors retain their immediate \(s\)-root, so their counts
do not change. Hence the tuple strictly decreases in lexicographic
order, shallow coordinate first, at every rewrite step.
There is no infinite descent in a fixed finite tuple of naturals:
the first coordinate can decrease only finitely often and then
stabilizes; repeat this argument inductively on the remaining
coordinates. This proves SN on each finite input, despite possible
growth in copied descendants and in total syntax size.

The embedding obstruction does not contradict this proof. Map
the left-side root to the new inner \(f\), its \(s\)-child
through the inserted \(p\) to the corresponding \(s\), and
all nodes of the captured \(x\) to that inner copy. Labels
and ancestry agree, with one inserted node on an edge. Reverse
embedding therefore prevents an orientation by a simplification
order of the earlier kind. The origin-count argument is a different
measure specialized to this actual one-rule system; no general
Puel tree-order theorem was used.
\end{proof}
\end{annexitem}

\begin{annexitem}{annex:p12-klop1987-pa-puel-extension-failure}{thm:p12-klop1987-pa-puel-extension-failure}{A companion rule that repeatedly reactivates the blocked f-node}
\begin{proof}
Take \(c\) to be an inert constant, or a fixed object variable
if there are no constants. In the enlarged ordinary system perform
\[
 f(s(c))\to g(s(c),f(p(s(c))))\to g(s(c),f(s(c))).
\]
The first step is \(Q\), and the second uses \(p(s(x))\to s(x)\)
inside the new \(f\)-argument. Put \(C[v]=g(s(c),v)\).
Context closure repeats the same two steps at the retained
\(f(s(c))\) inside \(C^n[f(s(c))]\), giving
\(C^{n+1}[f(s(c))]\) for every \(n\geq0\). Concatenating
these positive two-step paths gives an actual infinite reduction
of finite terms. The new \(f\) is no longer permanently inactive,
which is precisely the hypothesis that the preceding origin-count
proof used. Thus that proof cannot be promoted to arbitrary
systems containing the displayed source rule.
\end{proof}
\end{annexitem}
